\chapter{An introduction to censored, truncated or sample-selected data}

\section{Compare three different cases}


A sample is truncated if some observations are systematically excluded
from the sample.  For example, suppose we are interested in
relationship between people's income($y$) and education($x$).  If we
have observations of both $y$ and $x$ only for people whose income is
above \$20,000 per year, then we have a truncated sample.  A sample is
censored if no observations have been systematically excluded but some
of information contained in them has been suppressed.  In other words,
in a truncated sample, you don't have observations which have been
truncated; while in a censored sample, you do have observations but
you don't have full information.  The idea of "censoring" is that some
data above or below a threshold are mis-reported at the threshold.  If
we use the same example of people's income, if we don't observe
people's income ($y$) and education($x$) with income lower than
\$20,000 per year, then we have a truncated sample.  If we have
observations of $x$ for people whose income above AND below that
limit, but no information on exactly how much their income is for
those whose income below the limit (only thing we know is that it's
lower than \$20,000 per year), then we have a censored sample.



The problem of sample selection arises when no observations have been
systematically excluded but some information has been suppressed, just
like censoring.  However, in sample selection problems, we have
information on how the ``selection'' happens.  In censoring, usually a
constant ``threshold'' is set.  Observations on dependent variable for
those above (or below) that criteria are missing. Using the income
example again, the threshold is set to be \$20,000, we have no
information on exactly how much their income is for those whose income
below the limit (only thing we know is that it's lower than \$20,000
per year).  When we have a third variable that determines the
censoring situation, then we call it sample selection.  For example,
we have observed only union member's information on income, but we
observe information on education for everybody, including union
members or non-members.  In this example, union membership is an
additional variable that ``selects'' the sample which are involved in
the estimation of effect of education on income.

Richard Breen has a nice table in his book:


\begin{table}[htbp]\centering
 \caption{Comparison of censored, truncated, and sample-selected
 cases. (borrowed from Richard Breen)
\label{tabresult intreg}}
% Table generated by Excel2LaTeX from sheet 'Sheet1'
% Table generated by Excel2LaTeX from sheet 'Acr2891'


\begin{tabular}{rrrr}

   Sample  & Y Variable  & X Variable  &   Example  \\
\hline


 Censored  & y is known exactly  & x variables are  & Determinants of  \\

           & only if some  & observed for the  & income; income is  \\

           & criterion defined.  & entire sample,  & measured exactly  \\

           & in terms of y is  & regardless of  & only if it above the  \\

           &      met.  & whether y is  & poverty line. All  \\

           &            & observed exactly  & other incomes are  \\

           &            &            & reported at the  \\

           &            &            & poverty line  \\
 \hline

   Sample & y is observed only  & x and w (the  & Survey data with  \\

 Selected  & if a criteria defined.  & determinants of  & item or unit non- \\

           & in terms of some  & whether Z =1) are  &  response  \\

           & other random  & observed for the  &            \\

           & variable (Z) is met.  & entire sample,  &            \\

           &            & regardless of  &            \\

           &            & whether y is  &            \\

           &            & observed or not  &            \\

\hline

Truncated  & y is known only if  & x variables are  & Donations to  \\

           & some criterion  & observed only if y  & political  \\

           & defined in terms of  & is observed.  & campaigns.  \\

           & y is met.  &            &            \\

\end{tabular}
\end{table}



\section{An illustration for truncated or censored data}

Consider the model
\begin{equation}
y_t^0=\beta_1 +\beta_2 x_t +u_t, \ u_t \sim NID(0, \sigma^2),
\end{equation}
where $y_t^0$ is a latent variable.  We actually observe $y_t$, which
differs from $y_t^0$ , because it's either truncated or censored.
 
Suppose that censorship or truncation occurs whenever $y_t^0$ is less
than 0.  Clearly, the larger the error term $u_t$, the larger is
$y_t^0$ and thus the greater must be the probability that $y_t^0 \geq
0$.  This probability also depends on $x_t$.  So for the sample we
observe, $u_t$ does not have conditional mean 0 and is not
uncorrelated with $x_t$.  OLS using truncated or censored samples (OLS
of $y_t$ on $x_t$) yields biased and inconsistent estimators.
Normally we would like to draw inference on the population which is
represented by the full sample.  Shown in figure 1, ideally we would
have the OLS regression line (if we had the full sample), which is
very close to the ``true'' regression line (which is the mechnism
generates the data).  We would have the ``OLS censored sample'' line
if we run OLS on the censored sample; we would have ``OLS truncated
sample'' if we run OLS on the truncated sample.  It shows that both of
them are severely biased from the ``true'' regression line.

% \begin{figure}[htbp] \footnotesize \caption{\small Effects of truncation and censoring from Davidson and MacKinnon}
% \centerline{\includegraphics[width=4in,height=3in]{davidson.eps}}
% \label{fig:censor}
% \end{figure}

\section{Truncated Models} 

For truncated data, a consistent estimator comes from maximum
likelihood estimator (MLE).  If we assume the error terms in the
latent variable model has a known distribution, then MLE estimator can
be applied.  The most popular choice is Gaussian.
\begin{eqnarray}
\mbox{Pr}(y_t^0 \geq 0) & = & \mbox{Pr} ({\bf X}_t \beta +u_t \geq
0) \nonumber \\
& = & 1- \mbox{Pr} (u_t/\sigma < - {{\bf X}_t \beta} / \sigma)
\nonumber \\
 & = & 1- \Phi (-{{\bf X}_t \beta} / \sigma)=\Phi ({{\bf X}_t \beta} / \sigma)
\end{eqnarray}

The density of $y_t$ is proportional to the density of $y_t^0$
when $y_t^0 \geq 0$ and $y_t$ is observed.  It is 0 elsewhere. The
factor of proportionality, which is needed to ensure that the
density integrates to unity, is the inverse of the probability
that $y_t^0 \geq 0$.  The density of $y_t$ can be written as
\begin{equation}
\frac{\sigma^{-1} \phi ((y_t - {\bf X}_t \beta)/\sigma)}{\Phi
({\bf X}_t \beta / \sigma)}
\end{equation}

This implies the log-likelihood function,
\begin{equation}
\ell ({\bf y, \beta}, \sigma)=- \frac{n}{2} \log(2 \pi) -n
\log(\sigma) - \frac{1}{2 \sigma^2} \sum_{t=1}^n(y_t-{\bf X}_t
\beta)^2 - \sum_{t=1}^n \log \Phi ({\bf X}_t \beta / \sigma).
\end{equation}

This can be estimated by MLE.  In Stata, truncreg is the command to do
truncated regression model.

\section{Censored Models}

The most popular censored model is Tobit model.
\begin{eqnarray}
y_t^0 & = & {\bf X}_t +u_t, \quad u_t \sim NID(0, \sigma^2) \nonumber \\
y_t & = & y_t^0 \ \mbox{if} \ y_t^0 > 0; \ y_t=0 \
\mbox{otherwise}.
\end{eqnarray}

We see that
\begin{eqnarray}
\mbox{Pr}(y_t= 0) & = & \mbox{Pr} (y_t^0 \leq 0) = \mbox{Pr} ({\bf
X}_t \beta +u_t \leq 0) \nonumber \\
& = & \mbox{Pr} (u_t/\sigma < - {{\bf X}_t \beta} / \sigma)
\nonumber \\
 & = &  \Phi (-{{\bf X}_t \beta} / \sigma)
\end{eqnarray}

The contribution to the log-likelihood function made by
observations with $y_t=0$ is
\begin{equation}
\ell_t(y_t, \beta, \sigma) = \log \Phi (-{\bf X}_t \beta/\sigma).
\end{equation}

If $y_t$ is positive, the contribution to the log-likelihood is the
logarithm of the density,
\begin{equation}
\log(\frac{1}{\sigma} \phi ((y_t - {\bf X}_t \beta) / \sigma )).
\end{equation}

The log-likelihood function of the tobit model is
\begin{equation}
\sum_{y_t=0} \log \Phi (-{\bf X}_t \beta/\sigma) + \sum_{y_t>0}
\log(\frac{1}{\sigma} \phi ((y_t - {\bf X}_t \beta)/\sigma )))
\end{equation}

This can be estimated by MLE. In Stata, tobit or intreg can be used
for censoring models.

\section{Sample Selection}

The sample selection models differ from censored model in that it
involves a different variable (from $y$ itself) to determine the
censorship (selection).

Suppose that $y_t^0$ and $z_t^0$ are two latent variables,
generated by the bivariate process

\begin{equation}
\begin{bmatrix}
y_t^0 \\ z_t^0
\end{bmatrix}
=\begin{bmatrix} {\bf X}_t \beta \\ {\bf W}_t \gamma
\end{bmatrix}+
\begin{bmatrix} u_t \\ v_t
\end{bmatrix}, \
\begin{bmatrix} u_t \\ v_t
\end{bmatrix} \sim NID({\bf 0},
\begin{bmatrix} \sigma^2 & \rho  \\ \rho  & 1
\end{bmatrix})
\end{equation}

We observe $y_t$ and $z_t$:
\begin{eqnarray}
 y_t & = & y_t^0 \ \mbox{if} \ z_t^0 > 0;  \ y_t \ \mbox{unobservable otherwise}; \nonumber \\
 z_t & = & 1  \ \mbox{if} \ z_t^0 > 0; \ z_t=0 \ \mbox{otherwise}.
\end{eqnarray}

There are two types of observations, ones we observe $y_t=y_t^0$
and $z_t=1$, along with both $\bf X_t$ and $\bf W_t$, and ones we
observe only $z_t=0$ and $\bf W_t$.

Each observation contributes to the likelihood function by
\begin{equation}
I(z_t=0)\mbox{Pr}(z_t=0)+I(z_t=1)\mbox{Pr}(z_t=1)f(y_t^0 | z_t=1),
\end{equation}

Under the normality assumption,
\begin{equation}
u_t=\rho v_t+e_t
\end{equation}
where $e_t$ is independent of $v_t \sim N(0,1)$.  A useful fact
about the standard normal distribution is that
\begin{equation}
{\rm E} (v_t | v_t>-x)=\lambda (x) =\frac{\phi(x)}{\Phi(x)}
\end{equation}
and the function $\lambda(x)$ is called the inverse Mills ratio.

The log-likelihood function can be shown to be
\begin{equation}
\sum_{z_t=0} \log \Phi(-{\bf W}_t \gamma)+\sum_{z_t=1}
\log(\frac{1}{\sigma} \phi ((y_t - {\bf X}_t \beta)/\sigma)) +
\sum_{z_t=1} \log \Phi (\frac{{\bf W}_t \gamma+\rho (y_t-{\bf X}_t
\beta)/\sigma}{(1-\rho^2)^{1/2}})
\end{equation}

So this model can be estimated by MLE.

However, it is popular to use Heckman's two-step method.

Heckman's method is based on the fact that the original latent
model can be rewritten as
\begin{equation}
y_t={\bf X}_t \beta + \rho v_t +e_t
\end{equation}

Here the error term $u_t$ is divided into two parts, one
perfectly correlated with $v_t$, and one independent of $v_t$.

In the first step, an ordinary probit model is used to obtain
consistent estimates $\hat \gamma$ of the parameters of the
selection equation.

In the second step, the unobserved $v_t$ is replaced by the
selectivity regressor $\frac{\phi({\bf W}_t \hat
\gamma)}{\Phi({\bf W}_t \hat \gamma)}$.

Therefore, the regression becomes
\begin{equation}
y_t={\bf X}_t \beta + \rho \frac{\phi({\bf W}_t \hat
\gamma)}{\Phi({\bf W}_t \hat \gamma)} + e_t
\end{equation}

Or,
\begin{equation}
y_t={\bf X}_t \beta + \rho \hat \lambda_t + e_t
\end{equation}

One thing to note in Heckman model is that in many situations, $\hat
\lambda$ is believed to be highly collinear with $\bf X_t$ (See
Olsen), if $\bf W_t$ is the same as $\bf X_t$.  However, the model can
still be estimated due to the nonlinearity of the model.
Nevertheless, it becomes a regular practice to require $\bf W_t$
contains at least one extra variable than $\bf X_t$, which is
sometimes called exclusion restriction.  In many situations, $\bf W_t$
contains all $\bf X_t$ variables, and at least one more variable.  The
reason to contain all $\bf X_t$ variables is because the selection is
``endogenous'' in the sense that $y_t$ is a factor in determining
selection; this is modeled by including all the $\bf X_t$'s.


In Stata, heckman is the command to do sample selection models.  It
has options to do a Heckman two-step estimation or MLE estimation.
Usually MLE is preferred.

\section{Switching Regression (Treatment-Effects Model)}


There is another situation that is similar to sample selection model:
we observe $y$ for $z=1$ and $z=0$.  In the example of effect of
education on income, we observe both union member's income and
non-member's income.  In this case, we have switching regression model
or treatment-effect model.  The treatment effect model estimates the
effect of an endogenous binary treatment$z_t$ (treatment, program
participation, etc.) on a continuous, fully-observed variable,$y_t$,
conditional on the independent variables $x_t$ and $w_t$.

Suppose $z_t^0$ is a latent variable.

\begin{equation}
\begin{bmatrix}
y_t \\ z_t^0
\end{bmatrix}
=\begin{bmatrix} {\bf X}_t \beta \\ {\bf W}_t \gamma
\end{bmatrix}+ \delta \begin{bmatrix} {z_t} \\ 0
\end{bmatrix} +
\begin{bmatrix} u_t \\ v_t
\end{bmatrix}, \
\begin{bmatrix} u_t \\ v_t
\end{bmatrix} \sim NID({\bf 0},
\begin{bmatrix} \sigma^2 & \rho  \\ \rho  & 1
\end{bmatrix})
\end{equation}

We observe $y_t$ and $z_t$:
\begin{eqnarray}
  z_t  =  1  \ \mbox{if} \ z_t^0 > 0; \ z_t=0 \ \mbox{otherwise}.
\end{eqnarray}

Notice that the only difference between the switching regression and
selection model is that we observe $y_t$ when $z_t=0$ and $z_t=1$.
It's not a selection, but a regime switching.

Both Stata (treatreg) and SAS (qlim) can estimate the switching
regression.  The default is by MLE.
